% Teil II: der durchgerechnete Weg.

% Gedankenschritt: Die Eckenbilanz verknüpft lokale Grade mit globaler
% Knotenzahl; sie liefert die Blattzahl in einer Zeile.

\subsection*{Gegeben und gesucht}
Baum mit $n\ge2$ Knoten; jeder Knoten, der kein freies Ende ist, hat genau
$3$ Nachbarn. Zu zeigen: $n$ gerade und die Zahl $L$ der freien Enden ist
$\tfrac n2+1$.

\subsection*{Ansatz}
\textbf{Warum die Eckenbilanz?} Die Aufgabe kennt nur lokale Grade
(Blätter haben Grad 1, alle anderen Grad 3). Das Handschlaglemma ist die
einzige globale Aussage, die die Grade \emph{aller} Knoten mit der
Kantenzahl verbindet --- und in einem Baum ist die Kantenzahl $n-1$
bekannt. Damit haben wir ``eine Unbekannte ($L$) und eine Gleichung``.
Wir stellen sie auf und lösen.

\subsection*{Rechnung}
Sei $L$ die Zahl der freien Enden (Grad $1$), $n-L$ die der Grad-3-Knoten.
Der Baum hat $n-1$ Kanten, also
\begin{align}
  \underbrace{L\cdot1}_{=L} + \underbrace{(n-L)\cdot3}_{=3n-3L} &= 2(n-1), \\
  L+3n-3L &= 2n-2, \\
  3n-2L &= 2n-2, \\
  n+2 &= 2L,
  \qquad L=\frac{n+2}2=\frac n2+1 .
\end{align}
Da $L$ eine \emph{ganze Zahl} (Knotenzahl) ist, folgt aus $2L=n+2$:
$n=2L-2$ ist gerade. Damit sind beide Teile der Behauptung bewiesen.

\textbf{Der entscheidende Gedanke:} Es braucht gar keine Fallunterscheidung
oder Graph-Theorie --- die eine Bilanz über die Grade liefert die
Blattzahl \emph{direkt}. (Der alternative Induktionsbeweis kommt zu
demselben Ergebnis.)

\begin{ergebnis}
\begin{equation*}
  n\ \text{gerade},\qquad \text{Anzahl freier Enden } L=\frac n2+1 .
\end{equation*}
\end{ergebnis}

\begin{kontrolle}
\textbf{Kleine Bäume von Hand:}
\begin{itemize}
  \item $n=2$ (zwei Knoten, beide Blätter): $L=2=\frac22+1$. \checkmark
  \item $n=4$: ein Mittelknoten mit drei Blättern, $L=3=\frac42+1$.
        \checkmark
  \item $n=6$: $L=4$. Ein ``Y``-Muster plus Verlängerung: stimmt.
\end{itemize}

\textbf{Plausibilität:} Je größer $n$, desto mehr ``Verlängerungen`` einer
Kette, pro hinzugefügten inneren Knoten (der wieder ein Blatt trägt)
kommt ein neues Blatt dazu. Die Formel $L=\frac n2+1$ deckt sich mit
dieser bildhaften Zählung.

\textbf{Grenzfall:} Der kleinste Fall $n=2$ ist die entartete ``zwei
Knoten``-Struktur, ein einzelner Blattknoten $n=1$ wäre von der
Voraussetzung ``jeder Nicht-Blatt hat Grad 3`` gar nicht betroffen;
der sinnvolle Gültigkeitsbereich ist $n\ge2$ mit der konsistenten
Randlösung.
\end{kontrolle}

\begin{merksatz}
Gradsumme $=$ doppelte Kantenzahl macht aus lokalen Grad-Angaben eine
globale Aussage; das liefert $n$ gerade und die Blattzahl
$L=\frac n2+1$ in einem Schritt.
\end{merksatz}
